



\olpart{fol}{First-order Logic}

\begin{editorial}
  This part covers the metatheory of first-order logic through
  completeness.  Currently it does not rely on a separate treatment of
  propositional logic; everything is proved.  The source files will
  exclude the material on quantifiers (and replace ``!!{structure}''
  with ``!!{valuation}'', $\Struct{M}$ with $\pAssign{v}$, etc.) if
  the ``FOL'' tag is false. In fact, most of the material in the part
  on propositional logic is simply the first-order material with the
  ``FOL'' tag turned off.

  If the part on propositional logic is included, this results in a
  lot of repetition. It is planned, however, to make it possible to
  let this part take into account the material on propositional logic
  (and exclude the material already covered, as well as shorten proofs
  with references to the respective places in the propositional
  part).
\end{editorial}



\olchapter{fol}{int}{Introduction to First-Order Logic}



\olfileid{fol}{int}{fol}

\olsection{First-Order Logic}

You are probably familiar with first-order logic from your first
introduction to formal logic.\footnote{In fact, we more or less assume
you are!{} If you're not, you could review a more elementary textbook,
such as \emph{forall x} \citep{Magnus2021}.} You may know it as
``quantificational logic'' or ``predicate logic.''  First-order
logic, first of all, is a formal language.  That means, it has a
certain vocabulary, and its expressions are strings from this
vocabulary.  But not every string is permitted.  There are different
kinds of permitted expressions: terms, !!{formula}s, and
!!{sentence}s.  We are mainly interested in !!{sentence}s of
first-order logic: they provide us with a formal analogue of sentences
of English, and about them we can ask the questions a logician
typically is interested in. For instance: 
\begin{itemize}
    \item Does $!B$ follow from~$!A$ logically?
    \item Is $!A$ logically true, logically false, or
contingent?
    \item Are $!A$ and $!B$ equivalent?
\end{itemize}

These questions are primarily questions about the ``meaning'' of
!!{sentence}s of first-order logic.  For instance, a philosopher would
analyze the question of whether $!B$ follows logically from~$!A$ as
asking: is there a case where $!A$ is true but~$!B$ is false ($!B$
doesn't follow from~$!A$), or does every case that makes $!A$ true
also make~$!B$ true ($!B$ does follow from~$!A$)?  But we haven't been
told yet what a ``case'' is---that is the job of \emph{semantics}.  The
semantics of first-order logic provides a mathematically precise model
of the philosopher's intuitive idea of ``case,'' and also---and this
is important---of what it is for !!a{sentence}~$!A$ to be \emph{true
in} a case. We call the mathematically precise model that we will
develop !!a{structure}. The relation which makes ``true in'' precise,
is called the relation of \emph{satisfaction}.  So what we will define
is ``$!A$ is satisfied in~$\Struct{M}$'' (in symbols: $\Sat{M}{!A}$)
for !!{sentence}s~$!A$ and !!{structure}s~$\Struct{M}$. Once this is
done, we can also give precise definitions of the other semantical
terms such as ``follows from'' or ``is logically true.'' These
definitions will make it possible to settle, again with mathematical
precision, whether, e.g., $\lforall[x][(!A(x) \lif !B(x)),
\lexists[x][!A(x)] \Entails \lexists[x][!B(x)]]$. The answer will, of
course, be ``yes.'' If you've already been trained to symbolize
sentences of English in first-order logic, you will recognize this as,
e.g., the symbolizations of, say, ``All ants are insects, there are
ants, therefore there are insects.'' That is obviously a valid
argument, and so our mathematical model of ``follows from'' for our
formal language should give the same answer.

Another topic you probably remember from your first introduction to
formal logic is that there are \emph{!!{derivation}s}.  If you have
taken a first formal logic course, your instructor will have made you
practice finding such !!{derivation}s, perhaps even !!a{derivation}
that shows that the above entailment holds.  There are many different
ways to give !!{derivation}s: you may have done something called
``natural deduction'' or ``truth trees,'' but there are many others.
The purpose of !!{derivation} systems is to provide tools using which the
logicians' questions above can be answered: e.g., a natural deduction
!!{derivation} in which $\lforall[x][(!A(x) \lif !B(x))$ and
$\lexists[x][!A(x)]$ are premises and $\lexists[x][!B(x)]]$ is the
conclusion (last line) \emph{verifies} that $\lexists[x][!B(x)]$
logically follows from $\lforall[x][(!A(x) \lif !B(x))]$ and
$\lexists[x][!A(x)]$.  

But why is that?  On the face of it, !!{derivation} systems have nothing to do
with semantics: giving a formal !!{derivation} merely involves arranging symbols
in certain rule-governed ways; they don't mention ``cases'' or ``true
in'' at all.  The connection between !!{derivation} systems and semantics has
to be established by a meta-logical investigation. What's needed is a
mathematical proof, e.g., that a formal !!{derivation} of $\lexists[x][!B(x)]$
from premises $\lforall[x][(!A(x) \lif !B(x))]$ and
$\lexists[x][!A(x)]$ is possible, if, and only if, $\lforall[x][(!A(x)
\lif !B(x))$ and $\lexists[x][!A(x)]$ together
entail~$\lexists[x][!B(x)]]$.  Before this can be done, however, a
lot of painstaking work has to be carried out to get the definitions
of syntax and semantics correct.





\olfileid{fol}{int}{syn}

\olsection{Syntax}

We first must make precise what strings of symbols count as
!!{sentence}s of first-order logic.  We'll do this later; for now
we'll just proceed by example.  The basic building blocks---the
vocabulary---of first-order logic divides into two parts. The first
part is the symbols we use to say specific things or to pick out
specific things. We pick out things using !!{constant}s, and we say
stuff about the things we pick out using !!{predicate}s.  E.g, we
might use $\Obj a$ as !!a{constant} to pick out a single thing, and
then say something about it using the !!{sentence}~$\Atom{\Obj P}{\Obj
a}$.  If you have meanings for ``$\Obj a$'' and ``$\Obj P$'' in mind,
you can read $\Atom{\Obj P}{\Obj a}$ as a sentence of English (and you
probably have done so when you first learned formal logic).  Once you
have such simple !!{sentence}s of first-order logic, you can build
more complex ones using the second part of the vocabulary: the logical
symbols (connectives and quantifiers). So, for instance, we can form
expressions like $(\Atom{\Obj P}{\Obj a} \land \Atom{\Obj Q}{\Obj b})$
or~$\lexists[\Obj x][\Atom{\Obj P}{\Obj x}]$.

In order to provide the precise definitions of semantics and the rules
of our !!{derivation} systems required for rigorous meta-logical study, we
first of all have to give a precise definition of what counts as
!!a{sentence} of first-order logic. The basic idea is easy enough to
understand: there are some simple !!{sentence}s we can form from just
!!{predicate}s and !!{constant}s, such as~$\Atom{\Obj P}{\Obj a}$. And
then from these we form more complex ones using the connectives and
quantifiers. But what exactly are the rules by which we are allowed to
form more complex !!{sentence}s?  These must be specified, otherwise
we have not defined ``!!{sentence} of first-order logic'' precisely
enough. There are a few issues. The first one is to get the right
strings to count as !!{sentence}s. The second one is to do this in
such a way that we can give mathematical proofs about \emph{all}
!!{sentence}s. Finally, we'll have to also give precise definitions of
some rudimentary operations with !!{sentence}s, such as ``replace
every $\Obj x$ in~$!A$ by~$\Obj b$.'' The trouble is that the
quantifiers and !!{variable}s we have in first-order logic make it not
entirely obvious how this should be done. E.g., should $\lexists[\Obj
x][\Atom{\Obj P}{\Obj a}]$ count as !!a{sentence}? What about
$\lexists[\Obj x][\lexists[\Obj x][\Atom{\Obj P}{\Obj x}]]$? What
should the result of ``replace $\Obj x$ by~$\Obj b$ in $(\Atom{\Obj
P}{\Obj x} \land \lexists[\Obj x][\Atom{\Obj P}{\Obj x}])$'' be?





\olfileid{fol}{int}{fml}

\section{\usetoken{P}{formula}}

Here is the approach we will use to rigorously specify !!{sentence}s
of first-order logic and to deal with the issues arising from the use
of !!{variable}s. We first define a \emph{different} set of
expressions: !!{formula}s. Once we've done that, we can consider the
role !!{variable}s play in them---and on the basis of some other
ideas, namely those of ``free'' and ``bound'' !!{variable}s, we can
define what !!a{sentence} is (namely, !!a{formula} without free
!!{variable}s). We do this not just because it makes the definition of
``!!{sentence}'' more manageable, but also because it will be crucial
to the way we define the semantic notion of satisfaction.

Let's define ``!!{formula}'' for a simple first-order language, one
containing only a single !!{predicate}~$\Obj P$ and a single
!!{constant}~$\Obj a$, and only the logical symbols $\lnot$, $\land$,
and~$\lexists$. Our full definitions will be much more general:
we'll allow infinitely many !!{predicate}s and !!{constant}s. In fact,
we will also consider !!{function}s which can be combined with
!!{constant}s and !!{variable}s to form ``terms.'' For now, $\Obj a$
and the variables will be our only terms.  We do need infinitely many
!!{variable}s.  We'll officially use the symbols $\Obj v_0$, $\Obj
v_1$, \dots, as variables.

\begin{defn}
The set of \emph{!!{formula}s}~$\Frm$ is defined as follows:
\begin{enumerate}
\item\ollabel{fmls-atom} $\Atom{\Obj P}{\Obj a}$ and $\Atom{\Obj
  P}{\Obj v_i}$ are !!{formula}s ($i \in \Nat$).

\ollabel{fmls-not}If $!A$ is !!a{formula}, then $\lnot !A$ is
  !!{formula}.

If $!A$ and $!B$ are !!{formula}s, then $(!A \land
  !B)$ is !!a{formula}.

\ollabel{fmls-ex}If $!A$ is !!a{formula} and $x$ is !!a{variable},
  then $\lexists[x][!A]$ is !!a{formula}.

\ollabel{fmls-limit}Nothing else is !!a{formula}.
\end{enumerate}
\end{defn}

\olref{fmls-atom} tells us that $\Atom{\Obj P}{\Obj a}$ and
$\Atom{\Obj P}{\Obj v_i}$ are !!{formula}s, for any $i \in
\Nat$. These are the so-called \emph{atomic} !!{formula}s. They give
us something to start from.  The other clauses give us ways of forming
new !!{formula}s from ones we have already formed. So for instance, by
\olref{fmls-not}, we get that $\lnot \Atom{\Obj P}{\Obj v_2}$ is
!!a{formula}, since $\Atom{\Obj P}{\Obj v_2}$ is already !!a{formula}
by \olref{fmls-atom}. Then, by \olref{fmls-ex}, we get that
$\lexists[\Obj v_2][\lnot \Atom{\Obj P}{\Obj v_2}]$ is another
!!{formula}, and so on.  \olref{fmls-limit} tells us that \emph{only}
strings we can form in this way count as !!{formula}s. In particular,
$\lexists[\Obj v_0][\Atom{\Obj P}{\Obj a}]$ and $\lexists[\Obj
  v_0][\lexists[\Obj v_0][\Atom{\Obj P}{\Obj a}]]$ \emph{do} count as
!!{formula}s, and $(\lnot \Atom{\Obj P}{\Obj a})$ does not, because of
the extraneous outer parentheses.

This way of defining !!{formula}s is called an \emph{inductive
definition}, and it allows us to prove things about !!{formula}s using
a version of proof by induction called \emph{structural induction}.
These are discussed in a general way in \olref[mth][ind][idf]{sec} and
\olref[mth][ind][sti]{sec}, which you should review before delving
into the proofs later on. Basically, the idea is that if you want to
give a proof that something is true for all !!{formula}s, you show
first that it is true for the atomic !!{formula}s, and then that
\emph{if} it's true for any !!{formula}~$!A$ (and~$!B$), it's
\emph{also} true for $\lnot !A$, $(!A \land !B)$, and
$\lexists[x][!A]$. For instance, this proves that it's true for
$\lexists[\Obj v_2][\lnot \Atom{\Obj P}{\Obj v_2}]$: from the first
part you know that it's true for the atomic !!{formula}~$\Atom{\Obj
P}{\Obj v_2}$. Then you get that it's true for $\lnot \Atom{\Obj
P}{\Obj v_2}$ by the second part, and then again that it's true for
$\lexists[\Obj v_2][\lnot \Atom{\Obj P}{\Obj v_2}]$ itself. Since all
!!{formula}s are inductively generated from atomic !!{formula}s, this
works for any of them.





\olfileid{fol}{int}{sat}

\olsection{Satisfaction}

We can already skip ahead to the semantics of first-order logic once
we know what !!{formula}s are: here, the basic definition is that of 
!!a{structure}. For our simple language, !!a{structure}~$\Struct M$ has
just three components: a non-empty set $\Domain{M}$ called the
\emph{!!{domain}}, what $\Obj a$ picks out in~$\Struct M$, and what
$\Obj P$ is true of in~$\Struct M$.  The object picked out by~$\Obj a$
is denoted~$\Assign{\Obj a}{M}$ and the set of things $\Obj P$ is true
of by~$\Assign{\Obj P}{M}$. !!^a{structure}~$\Struct{M}$ consists of
just these three things: $\Domain{M}$, $\Assign{\Obj a}{M} \in
\Domain{M}$ and $\Assign{\Obj P}{M} \subseteq \Domain{M}$. The general
case will be more complicated, since there will be many !!{predicate}s
and !!{constant}s, the !!{constant}s can have more than one place, and
there will also be !!{function}s.

This is enough to give a definition of satisfaction for !!{formula}s
that don't contain !!{variable}s.  The idea is to give an inductive
definition that mirrors the way we have defined !!{formula}s. We
specify when an atomic formula is satisfied in~$\Struct{M}$, and then
when, e.g., $\lnot !A$ is satisfied in~$\Struct{M}$ on the basis of
whether or not $!A$ is satisfied in~$\Struct{M}$. E.g., we could
define:
\begin{enumerate}
  \item $\Atom{\Obj P}{\Obj a}$ is satisfied in~$\Struct{M}$ iff
  $\Assign{\Obj a}{M} \in \Assign{\Obj P}{M}$.
  \item $\lnot !A$ is satisfied in~$\Struct{M}$ iff $!A$ is not
  satisfied in~$\Struct{M}$.
  \item $(!A \land !B)$ is satisfied in~$\Struct{M}$ iff $!A$ is
  satisfied in~$\Struct{M}$, and $!B$ is satisfied in~$\Struct{M}$ as
  well.
\end{enumerate}
Let's say that $\Domain{M} = \{0, 1, 2\}$, $\Assign{\Obj a}{M} = 1$,
and $\Assign{\Obj P}{M} = \{1, 2\}$. This definition would tell us
that $\Atom{\Obj P}{\Obj a}$ is satisfied in~$\Struct{M}$ (since
$\Assign{\Obj a}{M} =  1 \in \{1,2\} = \Assign{\Obj P}{M}$). It tells
us further that $\lnot \Atom{\Obj P}{\Obj a}$ is not satisfied
in~$\Struct{M}$, and that in turn $\lnot\lnot \Atom{\Obj P}{\Obj
a}$ is and $(\lnot \Atom{\Obj P}{\Obj a} \land \Atom{\Obj P}{\Obj a})$
is not satisfied, and so on.

The trouble comes when we want to give a definition for the
quantifiers: we'd like to say something like, ``$\lexists[\Obj
v_0][\Atom{\Obj P}{\Obj v_0}]$ is satisfied iff $\Atom{\Obj P}{\Obj
v_0}$ is satisfied.'' But the !!{structure}~$\Struct{M}$ doesn't tell
us what to do about !!{variable}s. What we actually want to say is
that $\Atom{\Obj P}{\Obj v_0}$ is satisfied \emph{for some value
of~$\Obj v_0$}.  To make this precise we need a way to assign
!!{element}s of~$\Domain{M}$ not just to $\Obj a$ but also to~$\Obj
v_0$. To this end, we introduce !!{variable} \emph{assignments}.
!!^a{variable} assignment is simply a function~$s$ that maps
!!{variable}s to !!{element}s of~$\Domain{M}$ (in our example, to one
of $1$, $2$, or~$3$).  Since we don't know beforehand which
!!{variable}s might appear in !!a{formula} we can't limit which
!!{variable}s $s$ assigns values to. The simple solution is to
require that $s$ assigns values to \emph{all} !!{variable}s~$\Obj
v_0$, $\Obj v_1$, \dots\@ We'll just use only the ones we need.

Instead of defining satisfaction of !!{formula}s just relative to
!!a{structure}, we'll define it relative to
!!a{structure}~$\Struct{M}$ \emph{and} !!a{variable} assignment~$s$,
and write $\Sat{M}{!A}[s]$ for short. Our definition will now include
an additional clause to deal with atomic !!{formula}s containing
!!{variable}s:
\begin{enumerate}
  \item $\Sat{M}{\Atom{\Obj P}{\Obj a}}[s]$ iff
  $\Assign{\Obj a}{M} \in \Assign{\Obj P}{M}$.
  \item $\Sat{M}{\Atom{\Obj P}{\Obj v_i}}[s]$ iff
  $s(\Obj v_i) \in \Assign{\Obj P}{M}$.
  \item $\Sat{M}{\lnot !A}[s]$ iff not $\Sat{M}{!A}[s]$.
  \item $\Sat{M}{(!A \land !B)}[s]$ iff $\Sat{M}{!A}[s]$ and $\Sat{M}{!B}[s]$.
\end{enumerate}
Ok, this solves one problem: we can now say when $\Struct{M}$
satisfies $\Atom{\Obj P}{\Obj v_0}$ for the value~$s(\Obj v_0)$. To
get the definition right for $\lexists[\Obj v_0][\Atom{\Obj P}{\Obj
v_0}]$ we have to do one more thing:  We want to have that
$\Sat{M}{\lexists[\Obj v_0][\Atom{\Obj P}{\Obj v_0}]}[s]$ iff
$\Sat{M}{\Atom{\Obj P}{\Obj v_0}}[s']$ for \emph{some} way $s'$ of
assigning a value to~$\Obj v_0$. But the value assigned to~$\Obj v_0$
does not necessarily have to be the value that $s(\Obj v_0)$ picks
out.  We'll introduce a notation for that: if $m \in \Domain{M}$, then
we let $\Subst{s}{m}{\Obj v_0}$ be the assignment that is just
like~$s$ (for all !!{variable}s other than~$\Obj v_0$), except to
$\Obj v_0$ it assigns~$m$. Now our definition can be:
\begin{enumerate}\setcounter{enumi}{4}
  \item $\Sat{M}{\lexists[\Obj v_i][!A]}[s]$ iff
  $\Sat{M}{!A}[\Subst{s}{m}{\Obj v_i}]$ for some~$m \in \Domain{M}$.
\end{enumerate}
Does it work out? Let's say we let $s(\Obj v_i) = 0$ for all~$i \in
\Nat$. $\Sat{M}{\lexists[\Obj v_0][\Atom{\Obj P}{\Obj v_0}]}[s]$ iff
there is an $m \in \Domain{M}$ so that $\Sat{M}{\Atom{\Obj P}{\Obj
v_0}}[\Subst{s}{m}{\Obj v_0}]$. And there is: we can choose $m = 1$ or
$m = 2$. Note that this is true even if the value~$s(\Obj v_0)$
assigned to~$\Obj v_0$ by $s$ itself---in this case, $0$---doesn't do
the job. We have $\Sat{M}{\Atom{\Obj P}{\Obj v_0}}[\Subst{s}{1}{\Obj
v_0}]$ but not $\Sat{M}{\Atom{\Obj P}{\Obj v_0}}[s]$.

If this looks confusing and cumbersome: it is. But the added
complexity is required to give a precise, inductive definition of
satisfaction for all !!{formula}s, and we need something like it to
precisely define the semantic notions.  There are other ways of doing
it, but they are all equally (in)elegant.





\olfileid{fol}{int}{snt}

\olsection{\usetoken{P}{sentence}}

Ok, now we have a (sketch of a) definition of satisfaction (``true
in'') for !!{structure}s and !!{formula}s. But it needs this
additional bit---!!a{variable} assignment---and what we wanted is a
definition of !!{sentence}s. How do we get rid of assignments, and
what are !!{sentence}s?

You probably remember a discussion in your first introduction to
formal logic about the relation between !!{variable}s and quantifiers.
A quantifier is always followed by !!a{variable}, and then in the part
of the !!{sentence} to which that quantifier applies (its ``scope''),
we understand that the !!{variable} is ``bound'' by that quantifier.
In !!{formula}s it was not required that every !!{variable} has a
matching quantifier, and !!{variable}s without matching quantifiers
are ``free'' or ``unbound.''  We will take !!{sentence}s to be all
those !!{formula}s that have no free !!{variable}s.

Again, the intuitive idea of when an occurrence of !!a{variable} in
!!a{formula}~$!A$ is bound, which quantifier binds it, and when it is
free, is not difficult to get. You may have learned a method for
testing this, perhaps involving counting parentheses.  We have to
insist on a precise definition---and because we have defined
!!{formula}s by induction, we can give a definition of the free and
bound occurrences of !!a{variable}~$x$ in !!a{formula}~$!A$ also by
induction.  E.g., it might look like this for our simplified language:
\begin{enumerate}
  \item If $!A$ is atomic, all occurrences of $x$ in it are free (that
  is, the occurrence of $x$ in~$\Atom{\Obj P}{x}$ is free).
  \item If $!A$ is of the form $\lnot !B$, then an occurrence of~$x$
  in~$\lnot !B$ is free iff the corresponding occurrence of~$x$ is
  free in~$!B$ (that is, the free occurrences of variables in~$
  !B$ are exactly the corresponding occurrences in~$\lnot !B$).
  \item If $!A$ is of the form $(!B \land !C)$, then an occurrence of~$x$
  in~$(!B \land !C)$ is free iff the corresponding occurrence of~$x$ is
  free in~$!B$ or in~$!C$.
  \item If $!A$ is of the form $\lexists[x][!B]$, then no occurrence
  of $x$ in~$!A$ is free; if it is of the form $\lexists[y][!B]$ where
  $y$ is a different !!{variable} than~$x$, then an occurrence of~$x$
  in $\lexists[y][!B]$ is free iff the corresponding occurrence of~$x$
  is free in~$!B$.
\end{enumerate}

Once we have a precise definition of free and bound occurrences of
variables, we can simply say: !!a{sentence} is any !!{formula} without
free occurrences of !!{variable}s.





\olfileid{fol}{int}{sem}

\olsection{Semantic Notions}

We mentioned above that when we consider whether $\Sat{M}{!A}[s]$
holds, we (for convenience) let $s$ assign values to all !!{variable}s,
but only the values it assigns to !!{variable}s in~$!A$ are used.  In
fact, it's only the values of \emph{free} variables in~$!A$ that
matter. Of course, because we're careful, we are going to prove this
fact. Since !!{sentence}s have no free variables, $s$~doesn't matter
at all when it comes to whether or not they are satisfied in
!!a{structure}.  So, when $!A$ is !!a{sentence} we can define
$\Sat{M}{!A}$ to mean ``$\Sat{M}{!A}[s]$ for all~$s$,'' which as it
happens is true iff $\Sat{M}{!A}[s]$ for at least one~$s$. We need to
introduce !!{variable} assignments to get a working definition of
satisfaction for !!{formula}s, but for !!{sentence}s, satisfaction is
independent of the !!{variable} assignments.

Once we have a definition of ``$\Sat{M}{!A}$,'' we know what ``case''
and ``true in'' mean as far as !!{sentence}s of first-order logic are
concerned. On the basis of the definition of $\Sat{M}{!A}$ for
!!{sentence}s we can then define the basic semantic notions of
validity, entailment, and satisfiability.  A sentence is valid,
$\Entails !A$, if every !!{structure} satisfies it. It is entailed by
a set of !!{sentence}s, $\Gamma \Entails !A$, if every !!{structure}
that satisfies all the !!{sentence}s in~$\Gamma$ also satisfies~$!A$.
And a set of !!{sentence}s is satisfiable if some !!{structure}
satisfies all !!{sentence}s in it at the same time.

Because !!{formula}s are inductively defined, and satisfaction is in
turn defined by induction on the structure of !!{formula}s, we can use
induction to prove properties of our semantics and to relate the
semantic notions defined.  We'll collect and prove some of these
properties, partly because they are individually interesting, but
mainly because many of them will come in handy when we go on to
investigate the relation between semantics and !!{derivation} systems. In order
to do so, we'll also have to define (precisely, i.e., by induction)
some syntactic notions and operations we haven't mentioned yet.





\olfileid{fol}{int}{sub}

\olsection{Substitution}

We'll discuss an example to illustrate how things hang together, and
how the development of syntax and semantics lays the foundation for
our more advanced investigations later. Our !!{derivation} systems should let
us !!{derive} $\Atom{\Obj P}{\Obj a}$ from $\lforall[\Obj
v_0][\Atom{\Obj P}]{\Obj v_0}$. Maybe we even want to state this as a
rule of inference. However, to do so, we must be able to state it in
the most general terms: not just for $\Obj P$, $\Obj a$, and $\Obj
v_0$, but for any !!{formula}~$!A$, and term~$t$, and
!!{variable}~$x$. (Recall that !!{constant}s are terms, but we'll
consider also more complicated terms built from !!{constant}s and
!!{function}s.) So we want to be able to say something like,
``whenever you have !!{derive}d $\lforall[x][!A(x)]$ you are justified
in inferring~$!A(t)$---the result of removing $\lforall[x]$ and
replacing~$x$ by~$t$.'' But what exactly does ``replacing $x$ by~$t$''
mean? What is the relation between $!A(x)$ and~$!A(t)$?  Does this
always work?

To make this precise, we define the operation of \emph{substitution}.
Substitution is actually tricky, because we can't just replace
all~$x$'s in~$!A$ by~$t$, and not every~$t$ can be substituted for
any~$x$. We'll deal with this, again, using inductive definitions. But
once this is done, specifying an inference rule as ``infer $!A(t)$
from $\lforall[x][!A(x)]$'' becomes a precise definition.  Moreover,
we'll be able to show that this is a good inference rule in the sense
that $\lforall[x][!A(x)]$ entails~$!A(t)$. But to prove this, we have
to again prove something that may at first glance prompt you to ask
``why are we doing this?'' That $\lforall[x][!A(x)]$ entails~$!A(t)$
relies on the fact that whether or not $\Sat{M}{!A(t)}$ holds depends
only on the value of the term~$t$, i.e., if we let $m$ be whatever
!!{element} of~$\Domain{M}$ is picked out by~$t$, then
$\Sat{M}{!A(t)}[s]$ iff $\Sat{M}{!A(x)}[\Subst{s}{m}{x}]$. This holds
even when $t$ contains !!{variable}s, but we'll have to be careful
with how exactly we state the result.





\olfileid{fol}{int}{mod}

\olsection{Models and Theories}

Once we've defined the syntax and semantics of first-order logic, we
can get to work investigating the properties of !!{structure}s and the
semantic notions. We can also define !!{derivation} systems, and
investigate those.  For a set of !!{sentence}s, we can ask: what
!!{structure}s make all the !!{sentence}s in that set true?  Given a
set of !!{sentence}s~$\Gamma$, !!a{structure}~$\Struct{M}$ that
satisfies them is called a \emph{model of~$\Gamma$}.  We might start
from~$\Gamma$ and try to find its models---what do they look like? How
big or small do they have to be? But we might also start with a single
!!{structure} or collection of !!{structure}s and ask: what
!!{sentence}s are true in them?  Are there !!{sentence}s that
\emph{characterize} these !!{structure}s in the sense that they, and
only they, are true in them? These kinds of questions are the domain
of \emph{model theory}.  They also underlie the \emph{axiomatic
  method}: describing a collection of !!{structure}s by a set of
!!{sentence}s, the axioms of a theory. This is made possible by the
observation that exactly those !!{sentence}s entailed in first-order
logic by the axioms are true in all models of the axioms.

As a very simple example, consider preorders. A preorder is a
relation~$R$ on some set~$A$ which is both reflexive and transitive.
A set~$A$ with a two-place relation $R \subseteq A \times A$ on it is
exactly what we would need to give !!a{structure} for a first-order
language with a single two-place relation symbol~$\Obj P$: we would
set $\Domain{M} = A$ and $\Assign{\Obj P}{M} = R$.  Since $R$~is a
preorder, it is reflexive and transitive, and we can find a
set~$\Gamma$ of !!{sentence}s of first-order logic that say this:
\begin{align*}
  & \lforall[\Obj v_0][\Atom{\Obj P}{\Obj v_0,\Obj v_0}]\\
  & \lforall[\Obj v_0][\lforall[\Obj v_1][\lforall[\Obj v_2][((\Atom{\Obj P}{\Obj v_0,\Obj v_1} \land \Atom{\Obj P}{\Obj v_1,\Obj v_2}) \lif \Atom{\Obj P}{\Obj v_0,\Obj v_2})]]]
\end{align*}
These !!{sentence}s are just the symbolizations of ``for any~$x$,
$Rxx$'' ($R$ is reflexive) and ``whenever $Rxy$ and $Ryz$ then also
$Rxz$'' ($R$ is transitive). We see that !!a{structure}~$\Struct{M}$
is a model of these two !!{sentence}s~$\Gamma$ iff $R$ (i.e.,
$\Assign{\Obj P}{M}$), is a preorder on~$A$ (i.e., $\Domain{M}$). In
other words, the models of $\Gamma$ are exactly the preorders. Any
property of all preorders that can be expressed in the first-order
language with just~$\Obj P$ as !!{predicate} (like reflexivity and
transitivity above), is entailed by the two !!{sentence}s in~$\Gamma$
and vice versa.  So anything we can prove about models of~$\Gamma$ we
have proved about all preorders.

For any particular theory and class of models (such as $\Gamma$ and
all preorders), there will be interesting questions about what can be
expressed in the corresponding first-order language, and what cannot
be expressed. There are some properties of !!{structure}s that are
interesting for all languages and classes of models, namely those
concerning the size of the !!{domain}.  One can always express, for
instance, that the !!{domain} contains exactly $n$~!!{element}s, for
any~$n \in \PosInt$.  One can also express, using a set of infinitely
many !!{sentence}s, that the !!{domain} is infinite.  But one cannot
express that the domain is finite, or that the domain is
!!{nonenumerable}. These results about the limitations of first-order
languages are consequences of the compactness and L\"owenheim--Skolem
theorems.





\olfileid{fol}{int}{scp}

\olsection{Soundness and Completeness}

We'll also introduce !!{derivation} systems for first-order logic.  There are
many !!{derivation} systems that logicians have developed, but they all define
the same !!{derivability} relation between !!{sentence}s. We say that
$\Gamma$ \emph{!!{derive}s}~$!A$, $\Gamma \Proves !A$, if there is
!!a{derivation} of a certain precisely defined sort. !!^{derivation}s
are always finite arrangements of symbols---perhaps a list of
!!{sentence}s, or some more complicated structure.  The purpose of
!!{derivation} systems is to provide a tool to determine if !!a{sentence} is
entailed by some set~$\Gamma$.  In order to serve that purpose, it
must be true that $\Gamma \Entails !A$ if, and only if, $\Gamma
\Proves !A$.

If $\Gamma \Proves !A$ but not $\Gamma \Entails !A$, our !!{derivation} system
would be too strong, prove too much.  The property that if $\Gamma
\Proves !A$ then $\Gamma \Entails !A$ is called \emph{soundness}, and
it is a minimal requirement on any good !!{derivation} system. On the other
hand, if $\Gamma \Entails !A$ but not $\Gamma \Proves !A$, then our
!!{derivation} system is too weak, it doesn't prove enough.  The property that
if $\Gamma \Entails !A$ then $\Gamma \Proves !A$ is called
\emph{completeness}. Soundness is usually relatively easy to prove (by
induction on the structure of !!{derivation}s, which are inductively
defined). Completeness is harder to prove.

Soundness and completeness have a number of important consequences. If
a set of !!{sentence}s~$\Gamma$ !!{derive}s a contradiction (such as
$!A \land \lnot !A$) it is called \emph{inconsistent}. Inconsistent
$\Gamma$s cannot have any models, they are unsatisfiable. From
completeness the converse follows: any $\Gamma$ that is not
inconsistent---or, as we will say, \emph{consistent}---has a model. In
fact, this is equivalent to completeness, and is the form of
completeness we will actually prove.  It is a deep and perhaps
surprising result: just because you cannot prove $!A \land \lnot !A$
from $\Gamma$ guarantees that there is !!a{structure} that is as
$\Gamma$ describes it.  So completeness gives an answer to the
question: which sets of !!{sentence}s have models? Answer: all and only
consistent sets do.

The soundness and completeness theorems have two important
consequences: the compactness and the L\"owenheim--Skolem theorem.
These are important results in the theory of models, and can be used
to establish many interesting results. We've already mentioned two:
first-order logic cannot express that the !!{domain} of !!a{structure}
is finite or that it is !!{nonenumerable}.

Historically, all of this---how to define syntax and semantics of
first-order logic, how to define good !!{derivation} systems, how to prove that
they are sound and complete, getting clear about what can and cannot
be expressed in first-order languages---took a long time to figure out
and get right.  We now know how to do it, but going through all the
details can still be confusing and tedious.  But it's also important,
because the methods developed here for the formal language of
first-order logic are applied all over the place in logic, computer
science, and linguistics. So working through the details pays off in
the long run.



\OLEndChapterHook





